\documentclass[../main]{subfiles}
\begin{document}
\chapter{Introduction}

In today's world, there has been a rapid development in infrastructure, including buildings and bridges, to meet the needs of the 
people. Hence, the requirement for slender bridges is gaining importance to make it a sustainable construction by the optimal use 
of material. For spans greater than 150 m, the use of cable-stayed bridges, and for very high spans, the use of Suspension Bridges, 
is an efficient solution. One of the important criteria in most cable-stayed and Suspension bridges for optimization is their 
performance under dynamic wind loads, among which flutter can be highlighted. Flutter is the instability failure 
of structures under wind load at a critical velocity (of wind). It is caused due to the interaction of the bridge with the air, 
where coupling between the vertical and torsional vibration modes of the bridge which results in a rapid increase in amplitude 
(both vertical deflection and angular rotation) over time. This was the reason for the Tacoma Narrows Bridge's failure in 1940. 
Hence, finding this critical velocity and improving it is essential for safe and efficient design. A 3D bridge can be idealised 
as a 2D cross-section with 2 DOF system in vertical (Heave) and rotational (Pitch) direction. To find the critical velocity, 
the dynamic equilibrium equation must be solved with aerodynamic forces. The key challenge is quantifying the aerodynamic forces acting on the structure. In 1934, Theodorsen, a Norwegian American aerodynamicist, developed an 
analytical solution for the aerodynamic forces for unsteady flutter analysis of an airfoil \cite{TheodorsenGarrick1934}. These solutions depend on the reduced frequency, which quantifies the unsteadiness and time lag between the applied motion and the structure's response. However, since Theodorsen's theory was based on potential flow assumptions, it cannot be directly applied to bridges. In 1971, Scanlan reformulated this theory as a semi-empirical method using flutter derivatives \cite{scanlan1971} — shape constants that relate the aerodynamic forces linearly to the structure's velocity, angular rotation, and speed of angular rotation.

Flutter derivatives are functions of reduced velocity and cannot be directly computed; they require external CFD simulations or wind-tunnel tests. The approaches 
for finding the Flutter Derivatives can be broadly classified based on two criteria: the testing 
environment (CFD or Wind Tunnel) and the nature of the applied motion (Free vibration or forced 
motion). For the preliminary analysis of flutter, forced motion tests with CFD have greater 
importance due to their easier setup and lower cost. In CFD with forced motion tests, a single 
frequency method is usually used to find flutter derivatives, where a heave/pitch motion of one 
frequency is applied to the structure to find the Flutter derivative. Obtaining flutter derivatives with 
this approach is highly computationally costly because of the following reasons. Since flutter 
derivatives are functions of reduced velocity, multiple CFD simulations (by varying the frequency 
of input forced motion) must be performed to obtain them. For finding each point in a set of 8 
flutter derivatives, 2 CFD simulations are required. To obtain an accurate representation of each 
flutter derivative function, we need a minimum of 6 to 8 points on the flutter derivative curve. 
Hence, a total of 16 simulations is required. Each CFD simulation will take approximately 48 hours. 
Hence, to obtain a complete set of flutter derivatives, 768 hours of simulation time are required.  

\vspace{0.5 cm}

To overcome this, a multi-frequency vibration method was introduced by Lin Huang et al in 2011 \cite{huang2011}. 
By this method, instead of applying a single forced motion, we can superpose multiple forced 
motions of different frequencies and apply them as a forced motion in CFD. Then use mathematical 
tools, such as the Fourier Transform, to analyze the resultant aerodynamic force and obtain its 
individual components. In this way, the method can be theoretically n times as efficient (where n 
is the number of combinations). But according to the paper, the multi-frequency method reduced 
the simulation by only 18\%. The reason was to accurately identify the flutter derivative using a 
multi-frequency method, more simulation time is required than with a single-frequency method. 
After that, for a few years, there wasn't research focusing on this. In 2023, G. Martínez-López 
investigated the multi-frequency method to obtaining flutter derivatives via curve fitting with LES 
Turbulence modelling \cite{Martínez-López_2023}. From this work, it was identified that the selection of the frequency played 
an important role in the method's accuracy. The same result as with Lin Huang was observed. Ie, 
to obtain accurate Flutter derivatives using the multi-frequency method, the simulation time must 
be increased. This leaves several open questions:

\begin{flushleft}   
\textbf{1. Reason for this inaccuracy:} Why does the multi-frequency method not achieve the theoretical efficiency improvement?

\vspace{0.2cm}

\textbf{2. Effect of frequency selection:} How does the choice of frequencies affect accuracy, and how can frequencies be selected efficiently? Additionally, how much does frequency selection impact the simulation time and the significance of this inaccuracy with respect to the critical velocity?

\vspace{0.2cm}

\textbf{3. Other affecting parameters:} Are there additional parameters beyond frequency that influence accuracy, such as the amplitude of motion or the number of frequency combinations?
\end{flushleft}

The above open questions have been researched in this master's thesis in the following approach. 
If an inaccuracy in the multi-frequency method exists, then two possibilities of the source of error 
are there. First, with respect to CFD simulation, or due to the mathematical tool used for post-
processing. So, first signal analysis with FFT of arbitrary signals has been conducted. The FFT tool 
used in this work is from the SciPy module in Python. To investigate the source of error in CFD, a 
nonlinearity study was conducted on the cross-section. This is done because the multi-frequency 
method relies entirely on the assumption of a linear relationship between the applied motion and 
the resulting aerodynamic forces. Based on signal analysis, rules have been formulated for selecting 
the frequency for the multi-frequency method. Based on the Nonlinearity analysis, rules have been 
formulated for selecting the amplitude of the multi-frequency method. Based on the formulated 
rules, two sets of multifrequency CFD simulations have been conducted to verify the efficiency of 
the proposed framework.   
\vspace{0.5 cm}

\vspace{0.5 cm}

This document is structured as follows. Chapter 2 presents the theoretical background including an introduction to the multi-frequency method and the mesh convergence study. Chapter 3 focuses on signal analysis with the Fast Fourier Transform (FFT) to investigate the accuracy of frequency extraction. Chapter 4 examines the nonlinearity of the cross-section through CFD simulations to determine amplitude effects. Chapters 5 and 6 present two sets of multi-frequency CFD analyses validating the proposed framework. Finally, Chapter 7 synthesizes the findings and proposes a comprehensive framework for efficient use of the multi-frequency method with FFT.

\end{document}
